\section{Przegląd literatury i stan badań}
\label{sec:przeglad_literatury}

W literaturze przedmiotu paradygmat cyberbezpieczeństwa ulega fundamentalnej redefinicji. Konwencjonalne podejście \textit{detect and block} (oparte na sygnaturach antywirusowych i heurystyce) ustępuje miejsca doktrynie \textbf{Assume Breach} (założenie o nieuchronnym włamaniu), której wiodącą realizacją technologiczną jest sprzętowa kompartmentaryzacja stacji roboczej~\cite{rutkowska2010qubes,kapera2026methodology}.

\subsection{Kompartmentaryzacja a tradycyjne modele wirtualizacji}
Podstawy architektury Qubes OS sformułowali Rutkowska i Wojtczuk~\cite{rutkowska2010qubes}, wskazując, że monolityczne jądra współczesnych systemów operacyjnych (liczące miliony linii kodu) oraz podatne sterowniki sprzętowe uniemożliwiają zapewnienie integralności pojedynczej instancji systemu. W odróżnieniu od komercyjnych hipernadzorców typu 2 (takich jak VirtualBox czy VMware Workstation), instalowanych na wierzchu podatnego systemu operacyjnego gospodarza, Qubes OS implementuje hypervisor typu 1 (bare-metal) \textbf{Xen}~\cite{barham2003xen}.

W architekturze tej uprzywilejowana domena administracyjna (\texttt{dom0}) jest całkowicie pozbawiona łączności sieciowej. Sterowniki kart sieciowych i magistrali USB są wyizolowane w dedykowanych maszynach domenowych (\texttt{sys-net}, \texttt{sys-usb}), komunikujących się poprzez sprzętowy mechanizm IOMMU (Intel VT-d / AMD-Vi). Jak wykazali w 2026 roku Kapera i Niemiec~\cite{kapera2026methodology} w kompleksowych badaniach empirycznych opartych na macierzy MITRE ATT\&CK, architektura ta redukuje przestrzeń podatności stacji roboczej o ponad 80\% w porównaniu z wiodącym systemem Windows 11. Kluczowym wyróżnikiem jest uniemożliwienie trwałej infekcji (\textit{persistence}) oraz przemieszczania lateralnego (\textit{lateral movement}) poprzez zastosowanie maszyn efemerycznych (\textit{Disposable VMs} -- \texttt{DispVM}), których stan dyskowy jest bezpowrotnie niszczony w momencie zamknięcia procesu~\cite{kapera2026methodology}.

\subsection{Wyzwania operacyjne OSINT i higiena tożsamości cyfrowej}
Prowadzenie badań typu OSINT przez studentów i naukowców wiąże się z ryzykiem zaawansowanego profilowania cyfrowego. Zgodnie z badaniami Bazzella~\cite{bazzell2023osint}, nowoczesne platformy internetowe stosują pasywne i aktywne mechanizmy identyfikacji użytkownika, takie jak fingerprinting silnika renderującego (Canvas, WebGL), profilowanie stosu TCP/IP oraz analizę ciasteczek śledzących. 

Połączenie systemu Qubes OS z dedykowaną dystrybucją Whonix~\cite{whonix_design2024} realizuje rygorystyczny dwuetapowy model izolacji. Maszyna bramy sieciowej (\texttt{sys-whonix}) wymusza tunelowanie całego ruchu przez rozproszoną sieć Tor~\cite{dingledine2004tor}, podczas gdy stacja robocza (\texttt{anon-whonix}) nie posiada wiedzy o fizycznym adresie IP, MAC ani parametrach sprzętowych karty sieciowej. W rezultacie nawet potencjalny exploit przeglądarki klasy RCE (Remote Code Execution) nie pozwala atakującemu na deanonymizację badacza, gdyż nie ma on programistycznego dostępu do stosu sieciowego gospodarza.

\subsection{Bezpieczeństwo łańcucha dostaw w procesie programistycznym}
Dynamiczny rozwój otwartoźródłowych repozytoriów kodu doprowadził do gwałtownego wzrostu ataków na łańcuch dostaw oprogramowania (\textit{software supply chain attacks}). Ohm i wsp.~\cite{ohm2020backstabber} oraz Ladisa i wsp.~\cite{ladisa2023sok} sklasyfikowali wektory infekcji za pośrednictwem menedżerów pakietów (\texttt{npm}, \texttt{pip}, \texttt{crates.io}). Powszechną techniką jest \textit{typosquatting} oraz ukierunkowane wstrzykiwanie złośliwego kodu do skryptów instalacyjnych (np. \texttt{setup.py} lub hooków \texttt{preinstall}). 

Głównym celem takich ataków jest bezszelestna eksfiltracja poświadczeń zgromadzonych na stacji roboczej programisty: kluczy prywatnych SSH (\texttt{\textasciitilde/.ssh}), tokenów API do chmur publicznych oraz kluczy prywatnych GPG. W tradycyjnym systemie jednodomenowym dowolny uruchomiony proces posiada uprawnienia do odczytu całego katalogu domowego zalogowanego użytkownika. Zastosowanie kompartmentaryzacji w standardzie Split-SSH i Split-GPG w Qubes OS pozwala fizycznie odseparować magazyn kluczy kryptograficznych od zanieczyszczonego środowiska kompilacji.

\subsection{Ryzyka związane z automatyczną oceną prac studenckich}
Kwestia bezpieczeństwa weryfikacji prac laboratoryjnych z informatyki była badana m.in. przez Edwardsa i Shams~\cite{edwards2017grading}. Autorzy wykazali, że nadsyłane przez studentów zadania programistyczne często wykazują cechy kodu wrogiego -- zarówno w wyniku celowego działania (np. próby ucieczki z piaskownicy, odczytania szablonów oceniania), jak i rażących błędów implementacyjnych (wycieki pamięci, rekurencyjne pętle \texttt{fork()}, próby otwarcia tysięcy gniazd sieciowych).

Stosowane w uczelniach lekkie rozwiązania wirtualizacji kontenerowej (np. Docker) są powszechnie krytykowane w literaturze za brak twardej bariery izolacyjnej~\cite{sultan2019containers}. Kontenery współdzielą jądro z systemem gospodarza, przez co podatności jądra Linuksa (np. wyścigi uprawnień \textit{Dirty COW}, błędy w podsystemie \textit{io\_uring}) umożliwiają natychmiastowe przejęcie uprawnień administratora (\texttt{root}) na maszynie sprawdzającej.

Dotychczasowe badania nad infrastrukturą uczelnianą koncentrowały się na serwerowych klastrach wirtualizacji (np. zarządzanie środowiskami wirtualnymi i modelami językowymi w Proxmox~VE~\cite{zarnecki2026llm_proxmox}). W literaturze brakuje jednak opracowania, które adaptowałoby architekturę Qubes OS bezpośrednio do wieloaspektowego profilu akademickiej stacji roboczej, integrując w jednym spójnym środowisku bezpieczny OSINT, kompartmentaryzowany development oraz automatyczną detonację kodu studenckiego. Niniejsza praca wypełnia tę istotną lukę badawczą.
